%!TEX root = ../../thesis.tex
% generiert von thesis/scripts/make_appendix_prompts.py — nicht von Hand aendern
\subsection*{Phase 1: Business Understanding (generic tasks 1.1--1.4, block \texttt{bu})}
\begin{AppCPrompt}
1.1 Determine business objectives
The first objective of the data analyst is to thoroughly understand, from a
\AppCCont@^business perspective, what the client really wants to accomplish. Often the
\AppCCont@^client has many competing objectives and constraints that must be properly
\AppCCont@^balanced. The analyst’s goal is to uncover important factors, at the beginning,
\AppCCont@^that can influence the outcome of the project. A possible consequence of
\AppCCont@^neglecting this step is to expend a great deal of effort producing the right
\AppCCont@^answers to the wrong questions.
Background: Record the information that is known about the organization’s
\AppCCont@^business situation at the beginning of the project.
Business objectives: Describe the customer’s primary objective, from a business
\AppCCont@^perspective. In addition to the primary business objective, there are typically
\AppCCont@^other related business questions that the customer would like to address. For
\AppCCont@^example, the primary business goal might be to keep current customers by
\AppCCont@^predicting when they are prone to move to a competitor. Examples of related
\AppCCont@^business questions are “How does the primary channel (e.g., ATM, visit branch,
\AppCCont@^internet) a bank customer uses affect whether they stay or go?” or “Will lower
\AppCCont@^ATM fees significantly reduce the number of high-value customers who leave?”
Business success criteria: Describe the criteria for a successful or useful
\AppCCont@^outcome to the project from the business point of view. This might be quite
\AppCCont@^specific and able to be measured objectively, such as reduction of customer
\AppCCont@^churn to a certain level or general and subjective such as “give useful insights
\AppCCont@^into the relationships.” In the latter case it should be indicated who makes the
\AppCCont@^subjective judgment.
1.2 Assess situation
This task involves more detailed fact-finding about all of the resources,
\AppCCont@^constraints, assumptions and other factors that should be considered in
\AppCCont@^determining the data analysis goal and project plan. In the previous task, your
\AppCCont@^objective is to quickly get to the crux of the situation. Here, you want to
\AppCCont@^flesh out the details.
Inventory of resources: List the resources available to the project, including:
\AppCCont@^personnel (business experts, data experts, technical support, data mining
\AppCCont@^personnel), data (fixed extracts, access to live warehoused or operational
\AppCCont@^data), computing resources (hardware platforms) and software (data mining tools,
\AppCCont@^other relevant software).
Requirements, assumptions and constraints: List all requirements of the project
\AppCCont@^including schedule of completion, comprehensibility and quality of results and
\AppCCont@^security as well as legal issues. As part of this output, make sure that you are
\AppCCont@^allowed to use the data. List the assumptions made by the project. These may be
\AppCCont@^assumptions about the data that can be checked during data mining, but may also
\AppCCont@^include non-checkable assumptions about the business upon which the project
\AppCCont@^rests. It is particularly important to list the latter if they form conditions
\AppCCont@^on the validity of the results. List the constraints on the project. These may
\AppCCont@^be constraints on the availability of resources, but may also include
\AppCCont@^technological constraints such as the size of data that it is practical to use
\AppCCont@^for modeling.
Risks and contingencies: List the risks or events that might occur to delay the
\AppCCont@^project or cause it to fail. List the corresponding contingency plans; what
\AppCCont@^action will be taken if the risks happen.
Terminology: Compile a glossary of terminology relevant to the project. This may
\AppCCont@^include two components: (1) A glossary of relevant business terminology, which
\AppCCont@^forms part of the business understanding available to the project. Constructing
\AppCCont@^this glossary is a useful “knowledge elicitation” and education exercise. (2) A
\AppCCont@^glossary of data mining terminology, illustrated with examples relevant to the
\AppCCont@^business problem in question.
Costs and benefits: Construct a cost-benefit analysis for the project, which
\AppCCont@^compares the costs of the project with the potential benefit to the business if
\AppCCont@^it is successful. The comparison should be as specific as possible, for example
\AppCCont@^using monetary measures in a commercial situation.
1.3 Determine data mining goals
A business goal states objectives in business terminology. A data mining goal
\AppCCont@^states project objectives in technical terms. For example, the business goal
\AppCCont@^might be “Increase catalog sales to existing customers.” A data mining goal
\AppCCont@^might be “Predict how many widgets a customer will buy, given their purchases
\AppCCont@^over the past three years, demographic information (age, salary, city, etc.) and
\AppCCont@^the price of the item.”
Data mining goals: Describe the intended outputs of the project that enables the
\AppCCont@^achievement of the business objectives.
Data mining success criteria: Define the criteria for a successful outcome to
\AppCCont@^the project in technical terms, for example a certain level of predictive
\AppCCont@^accuracy or a propensity to purchase profile with a given degree of “lift.” As
\AppCCont@^with business success criteria, it may be necessary to describe these in
\AppCCont@^subjective terms, in which case the person or persons making the subjective
\AppCCont@^judgment should be identified.
1.4 Produce project plan
Describe the intended plan for achieving the data mining goals and thereby
\AppCCont@^achieving the business goals. The plan should specify the anticipated set of
\AppCCont@^steps to be performed during the rest of the project including an initial
\AppCCont@^selection of tools and techniques.
Project plan: List the stages to be executed in the project, together with
\AppCCont@^duration, resources required, inputs, outputs and dependencies. Where possible
\AppCCont@^make explicit the large-scale iterations in the data mining process, for example
\AppCCont@^repetitions of the modeling and evaluation phases. As part of the project plan,
\AppCCont@^it is also important to analyze dependencies between time schedule and risks.
\AppCCont@^Mark results of these analyses explicitly in the project plan, ideally with
\AppCCont@^actions and recommendations if the risks appear. Note: the project plan contains
\AppCCont@^detailed plans for each phase. For example, decide at this point which
\AppCCont@^evaluation strategy will be used in the evaluation phase. The project plan is a
\AppCCont@^dynamic document in the sense that at the end of each phase a review of progress
\AppCCont@^and achievements is necessary and an update of the project plan accordingly is
\AppCCont@^recommended. Specific review points for these reviews are part of the project
\AppCCont@^plan, too.
Initial assessment of tools and techniques: At the end of the first phase, the
\AppCCont@^project also performs an initial assessment of tools and techniques. Here, you
\AppCCont@^select a data mining tool that supports various methods for different stages of
\AppCCont@^the process, for example. It is important to assess tools and techniques early
\AppCCont@^in the process since the selection of tools and techniques possibly influences
\AppCCont@^the entire project.
\end{AppCPrompt}
\subsection*{Phase 2: Data Understanding (generic tasks 2.1--2.4, block \texttt{du})}
\begin{AppCPrompt}
2.1 Collect initial data
Acquire within the project the data (or access to the data) listed in the
\AppCCont@^project resources. This initial collection includes data loading if necessary
\AppCCont@^for data understanding. For example, if you apply a specific tool for data
\AppCCont@^understanding, it makes perfect sense to load your data into this tool. This
\AppCCont@^effort possibly leads to initial data preparation steps. Note: if you acquire
\AppCCont@^multiple data sources, integration is an additional issue, either here or in the
\AppCCont@^later data preparation phase.
Initial data collection report: List the dataset (or datasets) acquired,
\AppCCont@^together with their locations within the project, the methods used to acquire
\AppCCont@^them and any problems encountered. Record problems encountered and any solutions
\AppCCont@^achieved to aid with future replication of this project or with the execution of
\AppCCont@^similar future projects.
2.2 Describe data
Examine the “gross” or “surface” properties of the acquired data and report on
\AppCCont@^the results.
Data description report: Describe the data which has been acquired, including:
\AppCCont@^the format of the data, the quantity of data, for example number of records and
\AppCCont@^fields in each table, the identities of the fields and any other surface
\AppCCont@^features of the data which have been discovered. Does the data acquired satisfy
\AppCCont@^the relevant requirements?
2.3 Explore data
This task tackles the data mining questions, which can be addressed using
\AppCCont@^querying, visualization and reporting. These include: distribution of key
\AppCCont@^attributes, for example the target attribute of a prediction task; relations
\AppCCont@^between pairs or small numbers of attributes; results of simple aggregations;
\AppCCont@^properties of significant sub-populations; simple statistical analyses. These
\AppCCont@^analyses may address directly the data mining goals; they may also contribute to
\AppCCont@^or refine the data description and quality reports and feed into the
\AppCCont@^transformation and other data preparation needed for further analysis.
Data exploration report: Describe results of this task including first findings
\AppCCont@^or initial hypothesis and their impact on the remainder of the project. If
\AppCCont@^appropriate, include graphs and plots, which indicate data characteristics or
\AppCCont@^lead to interesting data subsets for further examination.
2.4 Verify data quality
Examine the quality of the data, addressing questions such as: is the data
\AppCCont@^complete (does it cover all the cases required)? Is it correct or does it
\AppCCont@^contain errors and if there are errors how common are they? Are there missing
\AppCCont@^values in the data? If so how are they represented, where do they occur and how
\AppCCont@^common are they?
Data quality report: List the results of the data quality verification; if
\AppCCont@^quality problems exist, list possible solutions. Solutions to data quality
\AppCCont@^problems generally depend heavily on both data and business knowledge.
\end{AppCPrompt}
\subsection*{Phase 3: Data Preparation (generic tasks 3.1--3.5, block \texttt{prep})}
\begin{AppCPrompt}
3.1 Select data
Decide on the data to be used for analysis. Criteria include relevance to the
\AppCCont@^data mining goals, quality and technical constraints such as limits on data
\AppCCont@^volume or data types. Note that data selection covers selection of attributes
\AppCCont@^(columns) as well as selection of records (rows) in a table.
Rationale for inclusion/exclusion: List the data to be included/excluded and the
\AppCCont@^reasons for these decisions.
3.2 Clean data
Raise the data quality to the level required by the selected analysis
\AppCCont@^techniques. This may involve selection of clean subsets of the data, the
\AppCCont@^insertion of suitable defaults or more ambitious techniques such as the
\AppCCont@^estimation of missing data by modeling.
Data cleaning report: Describe what decisions and actions were taken to address
\AppCCont@^the data quality problems reported during the verify data quality task of the
\AppCCont@^data understanding phase. Transformations of the data for cleaning purposes and
\AppCCont@^the possible impact on the analysis results should be considered.
3.3 Construct data
This task includes constructive data preparation operations such as the
\AppCCont@^production of derived attributes, entire new records or transformed values for
\AppCCont@^existing attributes.
Derived attributes: Derived attributes are new attributes that are constructed
\AppCCont@^from one or more existing attributes in the same record. Examples: area = length
\AppCCont@^* width.
Generated records: Describe the creation of completely new records. Example:
\AppCCont@^create records for customers who made no purchase during the past year. There
\AppCCont@^was no reason to have such records in the raw data, but for modeling purposes it
\AppCCont@^might make sense to explicitly represent the fact that certain customers made
\AppCCont@^zero purchases.
3.4 Integrate data
These are methods whereby information is combined from multiple tables or
\AppCCont@^records to create new records or values.
Merged data: Merging tables refers to joining together two or more tables that
\AppCCont@^have different information about the same objects. Example: a retail chain has
\AppCCont@^one table with information about each store’s general characteristics (e.g.,
\AppCCont@^floor space, type of mall), another table with summarized sales data (e.g.,
\AppCCont@^profit, percent change in sales from previous year) and another with information
\AppCCont@^about the demographics of the surrounding area. Each of these tables contains
\AppCCont@^one record for each store. These tables can be merged together into a new table
\AppCCont@^with one record for each store, combining fields from the source tables. Merged
\AppCCont@^data also covers aggregations. Aggregation refers to operations where new values
\AppCCont@^are computed by summarizing together information from multiple records and/or
\AppCCont@^tables. For example, converting a table of customer purchases where there is one
\AppCCont@^record for each purchase into a new table where there is one record for each
\AppCCont@^customer, with fields such as number of purchases, average purchase amount,
\AppCCont@^percent of orders charged to credit card, percent of items under promotion, etc.
3.5 Format data
Formatting transformations refer to primarily syntactic modifications made to
\AppCCont@^the data that do not change its meaning, but might be required by the modeling
\AppCCont@^tool.
Reformatted data: Some tools have requirements on the order of the attributes,
\AppCCont@^such as the first field being a unique identifier for each record or the last
\AppCCont@^field being the outcome field the model is to predict. It might be important to
\AppCCont@^change the order of the records in the dataset. Perhaps the modeling tool
\AppCCont@^requires that the records be sorted according to the value of the outcome
\AppCCont@^attribute. A common situation is that the records of the dataset are initially
\AppCCont@^ordered in some way but the modeling algorithm needs them to be in a fairly
\AppCCont@^random order. For example, when using neural networks it is generally best for
\AppCCont@^the records to be presented in a random order although some tools handle this
\AppCCont@^automatically without explicit user intervention. Additionally, there are purely
\AppCCont@^syntactic changes made to satisfy the requirements of the specific modeling
\AppCCont@^tool. Examples: removing commas from within text fields in comma-delimited data
\AppCCont@^files, trimming all values to a maximum of 32 characters.
\end{AppCPrompt}
\subsection*{Phase 4: Modeling (generic tasks 4.1--4.4, block \texttt{model})}
\begin{AppCPrompt}
4.1 Select modeling technique
As the first step in modeling, select the actual modeling technique that is to
\AppCCont@^be used. Whereas you possibly already selected a tool in business understanding,
\AppCCont@^this task refers to the specific modeling technique, e.g., decision tree
\AppCCont@^building with C4.5 or neural network generation with back propagation. If
\AppCCont@^multiple techniques are applied, perform this task for each technique
\AppCCont@^separately.
Modeling technique: Document the actual modeling technique that is to be used.
Modeling assumptions: Many modeling techniques make specific assumptions on the
\AppCCont@^data, e.g., all attributes have uniform distributions, no missing values
\AppCCont@^allowed, class attribute must be symbolic etc. Record any such assumptions made.
4.2 Generate test design
Before we actually build a model, we need to generate a procedure or mechanism
\AppCCont@^to test the model’s quality and validity. For example, in supervised data mining
\AppCCont@^tasks such as classification, it is common to use error rates as quality
\AppCCont@^measures for data mining models. Therefore, we typically separate the dataset
\AppCCont@^into train and test set, build the model on the train set and estimate its
\AppCCont@^quality on the separate test set.
Test design: Describe the intended plan for training, testing and evaluating the
\AppCCont@^models. A primary component of the plan is to decide how to divide the available
\AppCCont@^dataset into training data, test data and validation datasets.
4.3 Build model
Run the modeling tool on the prepared dataset to create one or more models.
Parameter settings: With any modeling tool, there are often a large number of
\AppCCont@^parameters that can be adjusted. List the parameters and their chosen value,
\AppCCont@^along with the rationale for the choice of parameter settings.
Models: These are the actual models produced by the modeling tool, not a report.
Model description: Describe the resultant model. Report on the interpretation of
\AppCCont@^the models and document any difficulties encountered with their meanings.
4.4 Assess model
The data mining engineer interprets the models according to his domain
\AppCCont@^knowledge, the data mining success criteria and the desired test design. This
\AppCCont@^task interferes with the subsequent evaluation phase. Whereas the data mining
\AppCCont@^engineer judges the success of the application of modeling and discovery
\AppCCont@^techniques more technically, he contacts business analysts and domain experts
\AppCCont@^later in order to discuss the data mining results in the business context.
\AppCCont@^Moreover, this task only considers models whereas the evaluation phase also
\AppCCont@^takes into account all other results that were produced in the course of the
\AppCCont@^project. The data mining engineer tries to rank the models. He assesses the
\AppCCont@^models according to the evaluation criteria. As far as possible he also takes
\AppCCont@^into account business objectives and business success criteria. In most data
\AppCCont@^mining projects, the data mining engineer applies a single technique more than
\AppCCont@^once or generates data mining results with different alternative techniques. In
\AppCCont@^this task, he also compares all results according to the evaluation criteria.
Model assessment: Summarize results of this task, list qualities of generated
\AppCCont@^models (e.g., in terms of accuracy) and rank their quality in relation to each
\AppCCont@^other.
Revised parameter settings: According to the model assessment, revise parameter
\AppCCont@^settings and tune them for the next run in the Build Model task. Iterate model
\AppCCont@^building and assessment until you strongly believe that you found the best
\AppCCont@^model(s). Document all such revisions and assessments.
\end{AppCPrompt}
\subsection*{Phase 5: Evaluation (generic tasks 5.1--5.3, block \texttt{eval})}
\begin{AppCPrompt}
5.1 Evaluate results
Previous evaluation steps dealt with factors such as the accuracy and generality
\AppCCont@^of the model. This step assesses the degree to which the model meets the
\AppCCont@^business objectives and seeks to determine if there is some business reason why
\AppCCont@^this model is deficient. Another option of evaluation is to test the model(s) on
\AppCCont@^test applications in the real application if time and budget constraints permit.
\AppCCont@^Moreover, evaluation also assesses other data mining results generated. Data
\AppCCont@^mining results cover models which are necessarily related to the original
\AppCCont@^business objectives and all other findings which are not necessarily related to
\AppCCont@^the original business objectives but might also unveil additional challenges,
\AppCCont@^information or hints for future directions.
Assessment of data mining results with respect to business success criteria:
\AppCCont@^Summarize assessment results in terms of business success criteria including a
\AppCCont@^final statement whether the project already meets the initial business
\AppCCont@^objectives.
Approved models: After model assessment with respect to business success
\AppCCont@^criteria, the generated models that meet the selected criteria become approved
\AppCCont@^models.
5.2 Review process
At this point the resultant model hopefully appears to be satisfactory and to
\AppCCont@^satisfy business needs. It is now appropriate to do a more thorough review of
\AppCCont@^the data mining engagement in order to determine if there is any important
\AppCCont@^factor or task that has somehow been overlooked. This review also covers quality
\AppCCont@^assurance issues, e.g., did we correctly build the model? Did we only use
\AppCCont@^attributes that we are allowed to use and that are available for future
\AppCCont@^analyses?
Review of process: Summarize the process review and highlight activities that
\AppCCont@^have been missed and/or should be repeated.
5.3 Determine next steps
According to the assessment results and the process review, the project decides
\AppCCont@^how to proceed at this stage. The project needs to decide whether to finish this
\AppCCont@^project and move on to deployment if appropriate or whether to initiate further
\AppCCont@^iterations or set up new data mining projects. This task includes analyses of
\AppCCont@^remaining resources and budget that influences the decisions.
List of possible actions: List the potential further actions along with the
\AppCCont@^reasons for and against each option.
Decision: Describe the decision as to how to proceed along with the rationale.
\end{AppCPrompt}
